\section{Efficiency and scale-factor measurement}\label{sec:method}
Let $O$ denote the offline event selection, $R$ the reference-trigger OR, and $S_c$ the signal-trigger OR for channel $c$. The measured efficiency in bin $i$ is
\begin{equation}\label{eq:eff}
\widehat\epsilon_{d/m,c,i}=\frac{N_{d/m,c,i}(O\cap R\cap S_c)}{N_{d/m,c,i}(O\cap R)},\qquad
\mathrm{SF}_{c,i}=\frac{\widehat\epsilon_{d,c,i}}{\widehat\epsilon_{m,c,i}}.
\end{equation}
The numerator is a subset of the denominator. Each event enters once even when multiple signal paths fire. Simulation counts represent the configured sum of non-trigger weights; the same weights are applied in the numerator and denominator. The conditional efficiency $P(S_c\mid O,R)$ represents $P(S_c\mid O)$ only when the residual reference-trigger bias is controlled.
The supplied one-dimensional figures already combine data and MC efficiencies with their SF ratio. These are used directly in the main results; selected numerator, denominator, and efficiency diagnostics are provided in Appendix~\ref{app:2024-support} and Appendix~\ref{app:2025-support-2025}. The leading/subleading convention is used in one-dimensional plots. For the final two-dimensional correction, same-flavour leptons are ordered by transverse momentum, whereas the mixed-flavour map uses $x=p_{\mathrm T}^{e}$ and $y=p_{\mathrm T}^{\mu}$.
The lowest occupied bins probe the turn-on region and should not be described as plateau bins without checking the offline thresholds and trigger response. Invalid, unpopulated, and inaccessible bins must not be assigned unity by default.\todo{Verify interval construction, MC weight treatment, exact bin edges, and underflow/overflow rules from code.}
